\chapter{Metodologia}
\label{chap:metodologia}

\section{Caracterização do Dataset}

Utilizou-se o \textit{Dataset of Multi-Modal Physiological Signals} 
(disponível em IEEE DataPort) \cite{naser2023multimodal}, contendo registros 
de 16 participantes expostos a estímulos visuais indutores de emoções. O 
dataset compreende quatro modalidades de biossinais: EEG (Emotiv EPOC+, 8 canais), 
ECG, EMG e fNIRS. Os sinais foram sincronizados e armazenados com marcadores 
temporais indicando os momentos de linha de base e estimulação.

Os 8 canais de EEG foram posicionados em AF7, AF8, F3, F4, PO7, PO8, PO3 e PO4, 
operando a uma frequência de amostragem de 512 Hz. As modalidades de ECG e EMG 
foram registradas com canais únicos a 250 Hz. O sistema de fNIRS registrou as 
concentrações de oxiemoglobina e desoxiemoglobina em dois canais a 16 Hz.

\section{Protocolo Experimental}

O protocolo de aquisição está subdividido em duas etapas principais. A etapa inicial compreende uma linha de base com duração de 30 segundos, durante a qual o participante permanece em estado de repouso para o estabelecimento dos parâmetros fisiológicos basais. A segunda etapa consiste em um estímulo emocional controlado com duração de 30 segundos, totalizando períodos de registro variáveis de aproximadamente 60 a 260 segundos por participante, dependendo da modalidade e do sujeito.

\section{Validação do Teorema de Nyquist}

As frequências de amostragem configuradas foram validadas conforme o teorema de amostragem de Nyquist. Para o sinal de EEG, com frequência máxima de 50 Hz e amostragem de 512 Hz, obteve-se uma razão de 5,12. No sinal de ECG, com frequência máxima de 40 Hz amostrado a 250 Hz, a razão calculada foi de 3,13. No EMG, considerando frequência de corte de 100 Hz e amostragem de 250 Hz, a razão foi de 1,25. Por fim, para o fNIRS, com frequência máxima de 2 Hz e amostragem de 16 Hz, a razão obtida foi de 4,00. Todos os parâmetros atendem plenamente ao critério de Nyquist, conforme detalhado na Tabela \ref{tab:nyquist_validacao}.

\begin{table}[htbp]
    \centering
    \caption{Validação das taxas de amostragem conforme o teorema de Nyquist}
    \label{tab:nyquist_validacao}
    \begin{tabular}{lcccc}
        \toprule
        \textbf{Modalidade} & $f_s$ (Hz) & $f_{max}$ (Hz) & $f_s / 2f_{max}$ & \textbf{Conforme} \\
        \midrule
        EEG & 512 & 50 & 5,12 & Sim \\
        ECG & 250 & 40 & 3,13 & Sim \\
        EMG & 250 & 100 & 1,25 & Sim \\
        fNIRS & 16 & 2 & 4,00 & Sim \\
        \bottomrule
    \end{tabular}
\end{table}

\section{Ambiente de Desenvolvimento}

O fluxo de processamento foi desenvolvido em linguagem Python, versão 3.13, integrando a biblioteca MNE-Python para o pré-processamento de biossinais. As bibliotecas NumPy e Pandas foram empregadas para manipulação de tensores e matrizes de dados, enquanto Matplotlib e SciPy foram utilizadas para visualização e processamento digital de sinais. O gerenciamento de dependências foi executado por meio da ferramenta uv e do arquivo de configuração pyproject.toml, assegurando a reprodutibilidade do ambiente computacional.

\section{Avaliação de Qualidade do Sinal (SQI)}
\label{sec:metodo_sqi}

No segundo estágio do pipeline, implementou-se um sistema de avaliação automática de qualidade (\textit{Signal Quality Index} --- SQI) para assegurar a validade fisiológica dos dados. O processo fundamenta-se na segmentação temporal e na extração de descritores estatísticos e espectrais.

\subsection{Segmentação de Janelas}
Os sinais brutos foram particionados em janelas de tempo fixas de $T = 5$ segundos. Esta duração foi selecionada para capturar características quase-estacionárias dos biossinais (especialmente EEG e ECG) enquanto permite uma resolução temporal adequada para a detecção de transientes de ruído e artefatos de movimento.

\subsection{Métricas de Qualidade Estatística}
Para cada janela $w_i$, calcularam-se métricas de momentos centrais:
\begin{itemize}
    \item \textbf{Curtose ($K$):} Utilizada para detectar \textit{outliers} e picos de tensão (\textit{spikes}) anormais, onde valores elevados indicam presença de ruído impulsivo ou transientes instrumentais.
    \item \textbf{Assimetria ($S$):} Medida da falta de simetria na distribuição de amplitudes, sensível a desvios de linha de base e saturação assimétrica (\textit{clipping}).
    \item \textbf{Amplitude IQR:} A amplitude interquartil foi empregada como descritor de dispersão robusto para identificar variações bruscas de magnitude associadas a movimentos do sujeito.
\end{itemize}

\subsection{Análise no Domínio da Frequência}
A integridade espectral foi avaliada por meio de:
\begin{itemize}
    \item \textbf{Relação Sinal-Ruído (SNR):} Estimada através da densidade espectral de potência ($PSD$) calculada pelo método de Welch. O SNR em dB foi definido pela razão entre a potência na banda de interesse fisiológico e a potência do ruído de fundo em bandas adjacentes.
    \item \textbf{Entropia Espectral ($H_{spec}$):} Calculada como a entropia de Shannon normalizada do espectro de potências. Sinais biológicos íntegros apresentam complexidade espectral característica, enquanto sinais contaminados por ruído branco ou interferência senoidal pura (60 Hz) apresentam desvios significativos nesta métrica.
\end{itemize}

\subsection{Classificação e Lógica de Rejeição}
A decisão de aceitação de um segmento baseou-se em um conjunto de regras heurísticas e limiares ($\theta$) específicos por modalidade $m$:
\begin{equation}
    f(w_i) = 
    \begin{cases} 
    \text{Aceito}, & \text{se } SNR > \theta_{SNR} \land K < \theta_{K} \land H > \theta_{H} \\
    \text{Rejeitado}, & \text{caso contrário}
    \end{cases}
\end{equation}
Adicionalmente, janelas com variância próxima a zero ($\sigma^2 < \epsilon$) foram automaticamente rotuladas como \textit{Loose Electrode} (Eletrodo Solto), enquanto variações de amplitude superiores a 10 vezes o desvio padrão global foram classificadas como \textit{Movement Artifact}.

\subsection{Calibração Empírica dos Limiares}
\label{sec:calibracao_limiares}

Os limiares de rejeição ($\theta$) foram estabelecidos por meio de uma abordagem 
empírica combinada com referências da literatura. O procedimento consistiu em:

\begin{enumerate}
    \item \textbf{Análise de Distribuição:} Calculou-se a distribuição estatística 
    das métricas (SNR, curtose, entropia) para todos os segmentos de todos os 
    sujeitos, identificando valores extremos e suas causas conhecidas.
    
    \item \textbf{Calibração por Modalidade:} Os limiares foram ajustados para cada 
    modalidade considerando:
    \begin{itemize}
        \item EEG: Sensibilidade alta a piscadas oculares e artefatos de movimento 
        (curtose elevada esperada).
        \item ECG: Presença de \textit{spikes} R (curtose muito alta esperada).
        \item EMG: Atividade muscular descontínua (entropia variável).
        \item fNIRS: Relação sinal-ruído intrinsecamente baixa.
    \end{itemize}
    
    \item \textbf{Validação Visual:} Os limiares foram validados por inspeção 
    visual de segmentos aceitos e rejeitados, conforme exemplificado na 
    Figura \ref{fig:sqi_comparison}.
\end{enumerate}

\section{Análise Estatística Descritiva e Testes de Hipótese}
\label{sec:metodo_estatistica}

No terceiro estágio, realizou-se uma caracterização estatística exaustiva dos biossinais para compreender as distribuições de amplitude e as dependências entre canais.

\subsection{Estatística Descritiva e Momentos de Alta Ordem}
Para cada canal de cada modalidade, computaram-se:
\begin{itemize}
    \item \textbf{Medidas de Tendência Central e Dispersão:} Média, mediana, variância, desvio padrão e intervalo interquartil (IQR).
    \item \textbf{Momentos de Alta Ordem:} Curtose e Assimetria (\textit{Skewness}), essenciais para caracterizar o desvio da normalidade e a presença de caudas pesadas no sinal.
\end{itemize}

\subsection{Testes de Normalidade e Homocedasticidade}
A natureza estocástica dos sinais foi avaliada por meio de:
\begin{itemize}
    \item \textbf{Teste de Shapiro-Wilk \cite{shapiro1965analysis}:} Aplicado como teste primário de normalidade (limitado a 5000 amostras).
    \item \textbf{Teste de Kolmogorov-Smirnov (K-S):} Utilizado como teste de aderência contra uma distribuição normal teórica.
    \item \textbf{Teste de Levene \cite{levene1960contributions}:} Empregado para avaliar a homocedasticidade (igualdade de variâncias) entre os canais de uma mesma modalidade, sendo robusto a desvios de normalidade.
\end{itemize}

\subsection{Análise de Correlação Inter-Modalidade}
Calculou-se a matriz de correlação de Pearson entre todos os canais sincronizados. Esta análise permite identificar redundâncias entre canais de EEG e acoplamentos fisiológicos entre diferentes modalidades (ex: influência da frequência cardíaca no sinal fNIRS).

A Tabela \ref{tab:sqi_threshold_justification} sintetiza a fundamentação dos 
limiares estabelecidos.

\begin{table}[htbp]
    \centering
    \caption{Justificativa empírica para os limiares de rejeição}
    \label{tab:sqi_threshold_justification}
    \begin{tabular}{p{2.5cm}p{3cm}p{3.5cm}}
        \toprule
        \textbf{Modalidade} & \textbf{Limiar} & \textbf{Fundamentação} \\
        \midrule
        EEG & SNR mín. $-5{,}0$ dB & Inferior ao ruído instrumental típico \\
        EEG & Curtose máx. $50{,}0$ & Captura piscadas sem rejeitar fisiológico \\
        ECG & Curtose máx. $100{,}0$ & Accomoda picos R pronunciados \\
        EMG & Entropia mín. $0{,}20$ & Detecta bursts musculares íntegros \\
        fNIRS & SNR mín. $-50{,}0$ dB & Ruído fotônico intrínseco \\
        \bottomrule
    \end{tabular}
\end{table}

\section{Limpeza e Correção de Sinais}
\label{sec:metodo_limpeza}

No quarto estágio do pipeline, implementou-se um conjunto de técnicas de pré-processamento para correção de artefatos e normalização estatística dos sinais biossignais.

\subsection{Filtragem no Domínio da Frequência}

A remoção de interferência elétrica foi realizada por meio de dois filtros digitais complementares:

\begin{itemize}
    \item \textbf{Filtro Notch (50 Hz):} Elimina a interferência da rede elétrica brasileira (60 Hz para EUA, 50 Hz para Europa).
    \item \textbf{Filtro Passa-Banda (Band-Pass):} Parâmetros específicos por modalidade para preservar a banda fisiológica de interesse:
    \begin{itemize}
        \item EEG: 0,5 -- 50 Hz (captura bandas delta a beta)
        \item ECG: 0,5 -- 50 Hz (preserva complexo QRS e ondas T)
        \item EMG: 20 -- 250 Hz (banda clássica de EMG de superfície)
        \item fNIRS: 0,01 -- 0,1 Hz (apenas variações hemodinâmicas lentas)
    \end{itemize}
\end{itemize}

\subsection{Interpolação de Lacunas}

Gaps temporais (ausências de sinal) foram identificados e preenchidos por interpolação linear para extensão máxima de 1 segundo. Lacunas maiores foram preservadas como segmentos inválidos para posterior rejeição.

\subsection{Winsorização}

A normalização de \textit{outliers} foi aplicada por Winsorização no percentil 5\textsuperscript{o} e 95\textsuperscript{o}, substituindo valores extremos pelos limiares correspondentes em vez de descartá-los. Esta abordagem preserva o número de amostras enquanto reduz a influência de artefatos impulsivos.

\subsection{Rejeição por Z-Score}

Após a Winsorização, outliers residuais foram identificados por z-scores superiores a 3,0 em valor absoluto, correspondendo a aproximadamente 0,27\% de rejeição sob distribuição normal teórica.

\section{Segmentação Temporal (Janelamento)}
\label{sec:metodo_segmentacao}

No quinto estágio do pipeline, implementou-se a segmentação dos sinais em janelas temporais de tamanho fixo, propagando as informações de qualidade (SQI) obtidas no Estágio 2 para garantir a integridade dos segmentos.

\subsection{Definição de Janelas de Tempo}

Os sinais limpos foram particionados em janelas de tempo fixas com duração de $T = 5$ segundos, selecionadas para capturar características fisiológicas estáveis enquanto mantêm resolução temporal adequada. A escolha de 5 segundos justifica-se por:
\begin{itemize}
    \item EEG: Captura pelo menos 2--3 ciclos completos de ondas alpha (8--12 Hz).
    \item ECG: Captura aproximadamente 5--6 batimentos cardíacos normais.
    \item EMG: Captura padrões de ativação muscular representativos.
    \item fNIRS: Permite análise de variações hemodinâmicas (banda 0,01--0,1 Hz).
\end{itemize}

\subsection{Janelas com Sobreposição}

Para maximizar a utilização dos dados, implementou-se suporte a janelas sobrepostas (\textit{overlapping windows}) com fração de sobreposição configurável. A taxa de sobreposição $o$ define a proporção do deslocamento entre janelas consecutivas:
\begin{equation}
    \text{stride} = T \times (1 - o)
\end{equation}
\subsection{Critérios de Estabilidade Intra-Janela}

Cada segmento foi submetido a testes de estacionariedade e homogeneidade para garantir a validade fisiológica:

\begin{itemize}
    \item \textbf{Coeeficiente de Variação (CV):} Garantiu-se que o CV da janela seja inferior a 30\%, indicando estabilidade de amplitude.
    \item \textbf{Teste de Dickey-Fuller Aumentado (ADF):} Avaliou-se a estacionariedade estatística do segmento, rejeitando janelas não-estacionárias.
\end{itemize}

\subsection{Propagação de Rejeição via SQI}

As informações de qualidade calculadas no Estágio 2 foram propagadas para o janelamento:
\begin{itemize}
    \item Segmentos rejeitados no SQI (Loose Electrode, Movement Artifact, SNR baixo) foram marcados como inválidos.
    \item Segmentos parcialmente válidos tiveram suas máscaras de rejeição armazenadas para exclusão em análises subsequentes.
\end{itemize}

\subsection{Métricas de Análise Inter-Janelas}

Para caracterizar a variabilidade entre segmentos de um mesmo sujeito, calculou-se:
\begin{itemize}
    \item \textbf{Variância Inter-Janelas:} Medida da dispersão de descritores estatísticos entre segmentos consecutivos.
    \item \textbf{Correlação Cruzada:} Análise de similaridade morfológica entre janelas adjacentes.
\end{itemize}

\subsection{Cache de Segmentos}

Os segmentos válidos foram armazenados em arquivos \texttt{.npz} (NumPy compressed archive) para uso eficiente nos estágios subsequentes de extração de atributos. Cada arquivo contém:
\begin{itemize}
    \item Sinais de todas as modalidades para o sujeito.
    \item Vetores de máscara indicando amostras válidas/inválidas.
    \item Metadados: timestamps, durações, limiares utilizados.
\end{itemize}

\section{Extração de Atributos}
\label{sec:metodo_extracao}

No sexto estágio do pipeline, descritores quantitativos são extraídos de cada janela temporal gerada no Estágio 5, construindo a matriz de atributos para aprendizado de máquina. A modalidade fNIRS foi definitivamente excluída neste estágio, confirmando o resultado do Estágio 5 em que menos de 2,5\% das janelas foram retidas.

\subsection{Atributos no Domínio do Tempo}

Para todas as modalidades ativas (EEG, ECG e EMG), os seguintes descritores temporais são calculados por canal e por janela:

\begin{itemize}
    \item \textbf{RMS} (\textit{Root Mean Square}): $\text{RMS} = \sqrt{\frac{1}{N}\sum_{i=1}^{N} x_i^2}$ — amplitude efetiva do sinal.
    \item \textbf{MAV} (\textit{Mean Absolute Value}): $\text{MAV} = \frac{1}{N}\sum_{i=1}^{N} |x_i|$ — indicador de ativação.
    \item \textbf{Variância:} $\sigma^2 = \frac{1}{N}\sum_{i=1}^{N}(x_i - \bar{x})^2$.
    \item \textbf{ZCR} (\textit{Zero Crossing Rate}): taxa de cruzamentos por zero, sensível à frequência dominante.
    \item \textbf{Parâmetros de Hjorth:} Atividade ($\sigma^2$), Mobilidade ($\sqrt{\sigma^2_{dx}/\sigma^2_x}$) e Complexidade (razão de mobilidades de primeira e segunda derivadas), que caracterizam espalhamento espectral sem transformada de Fourier.
    \item \textbf{Assimetria e Curtose:} momentos estatísticos de terceira e quarta ordem, respectivamente.
\end{itemize}

\subsection{Atributos no Domínio da Frequência}

A estimativa da Densidade Espectral de Potência (DEP) é realizada pelo método de Welch, com janela de Hann e \textit{nperseg} = $\min(256, N)$ amostras. Os atributos extraídos são:

\begin{itemize}
    \item \textbf{Potência Total:} integral da DEP pela regra trapezoidal.
    \item \textbf{Frequência Média e Mediana:} média ponderada pela potência e frequência onde a potência acumulada atinge 50\%.
    \item \textbf{Entropia Espectral:} $H = -\sum p_i \log_2 p_i$, onde $p_i$ é a DEP normalizada.
    \item \textbf{Potência por Bandas (EEG):} delta (0,5--4 Hz), theta (4--8 Hz), alpha (8--13 Hz), beta (13--30 Hz) e gamma (30--50 Hz).
\end{itemize}

\subsection{Atributos de Variabilidade Cardíaca (ECG)}

Para o ECG, atributos adicionais de Variabilidade da Frequência Cardíaca (VFC) são extraídos a partir da detecção de picos R por \textit{threshold} adaptativo com \texttt{scipy.signal.find\_peaks}:

\begin{itemize}
    \item \textbf{Domínio do tempo (HRV):} $\overline{RR}$, SDNN (desvio padrão dos intervalos RR), RMSSD (raiz quadrada da média dos quadrados das diferenças sucessivas) e pNN50 (proporção de diferenças sucessivas maiores que 50 ms).
    \item \textbf{Domínio da frequência (HRV):} potência nas bandas LF (0,04--0,15 Hz) e HF (0,15--0,40 Hz), e razão LF/HF, estimadas sobre o tacograma RR reamostrado a 4 Hz.
\end{itemize}

\section{Estrutura do Pipeline de Processamento}

O pipeline projetado compreende dez estágios sequenciais. O estágio inicial corresponde à aquisição e validação de dados brutos. O segundo estágio contempla o cálculo de índices de qualidade (SQI) para a rejeição de segmentos. O terceiro estágio executa análises estatísticas e testes de hipótese. O quarto estágio aplica técnicas de filtragem e correção de artefatos. O quinto estágio realiza a segmentação temporal dos sinais em janelas. O sexto estágio executa a extração de atributos nos domínios do tempo e da frequência. O sétimo estágio realiza a engenharia de novos atributos. O oitavo estágio aplica redução de dimensionalidade por componentes principais. O nono estágio executa a seleção de variáveis informativas. O décimo estágio valida estatisticamente o conjunto final. Este relatório concentra-se na execução dos cinco primeiros estágios do projeto.
